\chapter{Interview Question Bank}
\index{interview questions}\index{question bank}%
This appendix consolidates the Interview Q\&A sections of Chapters~1--24 into a single
bank, organized by domain and tagged by difficulty. Answers are abbreviated to the core
argument; the chapter reference holds the full explanation, diagrams, and labs. Use it for
spaced repetition, mock interviews, and capstone-defense rehearsal.

\begin{notebox}
\textbf{Difficulty tags.} \textbf{[F]} foundational---expected of any candidate;
\textbf{[I]} intermediate---separates practitioners from readers; \textbf{[A]}
advanced---staff-level trade-off and failure-mode reasoning. Answer aloud in under ninety
seconds; if you cannot, rebuild the referenced chapter's lab, not just its prose.
\end{notebox}

\section{Platform and Workspace (Chapter 1)}
\index{interview questions!platform}%
\begin{itemize}
    \item \textbf{[F] Why does Databricks separate the control plane from the data plane?}
    (Ch.~1) The SaaS control plane orchestrates while data and compute stay in the customer's
    cloud account---data residency, network isolation, existing cloud commitments.
    \item \textbf{[F] All-purpose or job cluster for a nightly pipeline?} (Ch.~1) Job cluster:
    created per run, cheaper, workload-exact, destroyed after---no idle spend or shared state.
    \item \textbf{[I] What problem do cluster policies actually solve?} (Ch.~1) They make
    standards enforceable: budgets, node restrictions, and required tags become properties of
    every cluster, not documentation.
    \item \textbf{[I] How would you give 200 users compute without 200 bill shocks?} (Ch.~1)
    Policy-bound clusters per team, serverless where possible, job clusters for schedules,
    weekly per-tag chargeback from \texttt{system.billing.usage}.
    \item \textbf{[A] What would make you split into a second workspace?} (Ch.~1) Only a hard
    trust or environment boundary; performance and organization are handled by catalogs and
    policies.
\end{itemize}

\section{Delta Lake and Medallion Design (Chapters 2, 3, 5)}
\index{interview questions!Delta Lake}%
\begin{itemize}
    \item \textbf{[F] How can a file store be ACID?} (Ch.~2) Files are immutable; the atomic
    unit is the single-file commit to the transaction log---readers see a version only after
    its commit exists.
    \item \textbf{[I] Two jobs MERGE into the same table simultaneously---what happens?}
    (Ch.~2) Optimistic concurrency: one commits first; the other re-reads and reconciles
    (disjoint keys) or retries/fails on conflict.
    \item \textbf{[I] \texttt{RESTORE} versus re-running the pipeline to fix a bad load?}
    (Ch.~2) RESTORE for table-level mistakes (instant, auditable); re-run when the bug is in
    the transformation logic itself.
    \item \textbf{[I] Why not set VACUUM retention to 1 hour?} (Ch.~2) It deletes the files
    time travel and long-running readers need; retention must exceed the recovery objective
    and the longest query.
    \item \textbf{[A] Does Delta's ACID protect against a logically wrong write?} (Ch.~2)
    No---ACID guarantees atomicity and durability, not correctness. Correctness comes from
    constraints, expectations, and reconciliation.
    \item \textbf{[F] Why keep Bronze forever when storage costs money?} (Ch.~3) Bronze is the
    only answer to ``what did the source actually say'' and the only input for rebuilding
    after logic changes.
    \item \textbf{[I] How do you prove a transition lost no rows?} (Ch.~3) Reconciliation
    arithmetic per run: input $=$ accepted $+$ quarantined, asserted and logged to a run
    ledger.
    \item \textbf{[I] MERGE or \texttt{replaceWhere} for Silver?} (Ch.~3) MERGE for keyed
    upserts; \texttt{replaceWhere} for full-partition restatements. Both idempotent; never
    mixed per table.
    \item \textbf{[A] Why does the medallion make streaming easier later?} (Ch.~3) Each layer
    already has idempotent, keyed writes---exactly what streaming sinks require. Batch becomes
    streaming by changing triggers, not contracts.
    \item \textbf{[F] Why do small files hurt even at the same total bytes?} (Ch.~5) Per-file
    costs (listing, footers, scheduling) scale with file count; metadata work dominates and
    the driver becomes the bottleneck.
    \item \textbf{[I] Z-ORDER versus liquid clustering?} (Ch.~5) Both co-locate rows for
    skipping; liquid keeps keys as changeable metadata maintained incrementally. New tables:
    liquid; legacy: Z-ORDER until migrated.
    \item \textbf{[A] How would you convince a skeptic that liquid clustering helped?}
    (Ch.~5) Same query, same data, controlled medians across cold and warm runs, plus the
    file-statistics story. Feelings are not benchmarks.
\end{itemize}

\section{Spark Execution and Performance (Chapters 4, 9)}
\index{interview questions!Spark}%
\begin{itemize}
    \item \textbf{[F] Walk through what happens on \texttt{df.write.saveAsTable()}.} (Ch.~4)
    Catalyst optimizes the logical plan, the physical plan splits into stages at shuffles,
    one task per partition runs on executors, and the final commit publishes the table
    version.
    \item \textbf{[I] Why is the shuffle the usual suspect?} (Ch.~4) It is the only all-to-all
    data movement---network, serialization, spill---plus a barrier. Most big-job pathologies
    are shuffle pathologies.
    \item \textbf{[I] What does AQE give you, and when do you override it?} (Ch.~4) Runtime
    coalescing, broadcast conversion, and skew splitting from actual statistics; override on
    known-skewed keys or a provably fittable broadcast it declined.
    \item \textbf{[I] When do you cache?} (Ch.~4) When one expensive DataFrame feeds multiple
    actions. Never for single-scan pipelines; always unpersist on shared clusters.
    \item \textbf{[A] How do you size a cluster for a brand-new job?} (Ch.~4) Estimate shuffle
    volume, partition for $\sim$100--200\,MB per task, memory for the largest task plus spill
    headroom---then let three runs of UI metrics correct the estimate.
    \item \textbf{[I] When is salting better than AQE skew handling?} (Ch.~9) Extreme, known
    skew (one dominant key): salting spreads it by construction; AQE splits only after the
    plan exists.
    \item \textbf{[F] Why are pandas UDFs faster than Python UDFs?} (Ch.~9) Arrow columnar
    batches cross the JVM/Python boundary once per batch, not once per row.
    \item \textbf{[A] When is a job ``tuned enough''?} (Ch.~9) When the next measured
    improvement costs more engineering time than the compute it saves, given the SLA
    margin---tuning has an economic stopping rule.
\end{itemize}

\section{SQL Warehousing and BI (Chapters 6, 7)}
\index{interview questions!SQL and BI}%
\begin{itemize}
    \item \textbf{[F] Why is Databricks SQL different from a BI tool on a Spark cluster?}
    (Ch.~6) Photon execution, warehouse-level caching and scaling, governance integration,
    per-query profiles---a serving layer, not a general cluster with a SQL cell.
    \item \textbf{[I] Materialized view or scheduled Gold table?} (Ch.~6) MV for stable
    derived shapes with freshness tolerance; scheduled Gold when the logic is versioned
    business logic needing tests and review.
    \item \textbf{[I] A dashboard is slow only at 9 a.m. Diagnose.} (Ch.~6) Queueing at the
    concurrency bound: raise max clusters for the burst, move the dashboard to an MV, confirm
    serverless for seconds-fast scale-up.
    \item \textbf{[I] What does the query profile tell you that duration does not?} (Ch.~6)
    Where the time went: which operator scanned what and spilled what. Duration says slow; the
    profile says why.
    \item \textbf{[I] How do you stop Genie from hallucinating metrics?} (Ch.~7) Scope to
    certified tables, encode definitions as instructions, provide example queries, and measure
    first-try accuracy on a fixed question set.
    \item \textbf{[A] Who should a published dashboard run as?} (Ch.~7) A service identity
    whose access equals what subscribers may see; publishing as a privileged user silently
    escalates every subscriber's access.
\end{itemize}

\section{Governance and Security (Chapters 8, 21, 22)}
\index{interview questions!governance}\index{interview questions!security}%
\begin{itemize}
    \item \textbf{[F] Managed or external table---how do you decide?} (Ch.~8) Managed by
    default; external only when a non-Databricks writer owns the files or the location is
    mandated.
    \item \textbf{[I] Where do you grant---catalog, schema, or table?} (Ch.~8) The highest
    level matching intent: schema for ``analysts read Gold,'' table only for true exceptions.
    High-level grants are reviewable; a thousand table grants are not.
    \item \textbf{[I] Row filters/column masks versus dynamic views?} (Ch.~8) Filters and
    masks live on the table and apply to every consumer---prefer them; dynamic views for
    compound or shared logic.
    \item \textbf{[A] How do you prove least privilege to an auditor?} (Ch.~8) The grant
    matrix, a negative test per principal class, and \texttt{system.access.audit}
    samples---implemented, tested, logged.
    \item \textbf{[A] Why test grants as a restricted principal rather than reviewing GRANTs?}
    (Ch.~8) Grants compose (ownership, inheritance, nesting) in ways reviewers misread; only
    the negative test survives composition.
    \item \textbf{[F] Why must production jobs run as service principals?} (Ch.~21) Person-run
    jobs break when the person leaves, inherit excessive access, and blur accountability;
    principals are scoped, owned, rotated, auditable.
    \item \textbf{[I] Is an IP access list sufficient network security?} (Ch.~21) No---it is a
    perimeter filter; stolen credentials from an allowed network pass. Layer with SSO/MFA,
    short-lived tokens, private connectivity, and audit alerts.
    \item \textbf{[A] Which single control gives the most risk reduction per hour?} (Ch.~21)
    Running production workloads as least-privilege service principals: caps blast radius,
    survives personnel change, makes actions attributable.
    \item \textbf{[I] What does column-level lineage change operationally?} (Ch.~22) Impact
    analysis becomes a query---every downstream consumer enumerable before a change, every
    metric traceable after an anomaly.
    \item \textbf{[A] How do you handle a data-subject erasure request in a lakehouse?}
    (Ch.~22) Targeted deletes through the layers with deletion vectors, a provenance query for
    every copy, \texttt{VACUUM} timing understood---erasure is a pipeline with evidence.
\end{itemize}

\section{Ingestion, Pipelines, and Orchestration (Chapters 10--12)}
\index{interview questions!pipelines}%
\begin{itemize}
    \item \textbf{[F] What does DLT do that a notebook chain does not?} (Ch.~10) Dependency
    inference, execution-mode management, checkpointing, retries, per-dataset quality metrics,
    an event log, and UC lineage---the operational layer you would otherwise hand-build.
    \item \textbf{[I] Streaming table versus materialized view?} (Ch.~10) A streaming table
    appends each source record exactly once (facts); a materialized view maintains a query
    result (derived state). ``What arrived'' versus ``what is true.''
    \item \textbf{[I] \texttt{expect\_or\_drop} versus \texttt{expect\_or\_fail}?} (Ch.~10)
    Drop when the row is expendable and the pipeline must continue (with quarantine
    visibility); fail when bad data invalidates the whole batch.
    \item \textbf{[I] How does Auto Loader achieve exactly-once?} (Ch.~11) The checkpoint
    durably records processed files per micro-batch; restarts resume from it, and idempotent
    Delta sink commits complete the contract.
    \item \textbf{[A] What does \texttt{SEQUENCE BY} protect against?} (Ch.~11) Out-of-order
    and replayed CDC events: without it, an older update overwrites a newer one and silently
    regresses table state.
    \item \textbf{[I] What belongs in the ingestion contract with an external producer?}
    (Ch.~11) Schema and compatibility rules, key and sequence semantics, arrival SLA, and the
    malformed-data policy.
    \item \textbf{[I] What makes a repair run safe?} (Ch.~12) Task idempotency plus true
    dependency modeling: repair reuses successful outputs and re-runs only the failed subtree.
    \item \textbf{[A] How do you alert on trajectory, not breach?} (Ch.~12) Track per-task
    duration baselines; compare elapsed $+$ remaining-p50 against the SLA and page early.
    Breach alerts tell you what you already lost.
\end{itemize}

\section{Streaming and Real-Time (Chapters 13--16)}
\index{interview questions!streaming}%
\begin{itemize}
    \item \textbf{[F] Explain micro-batch execution to a backend engineer.} (Ch.~13) The
    stream is an infinite table; each trigger runs a small incremental batch over rows since
    the last committed offset, writes transactionally, and commits offsets to the checkpoint.
    \item \textbf{[I] What exactly does the checkpoint make safe?} (Ch.~13) Restart position,
    stateful-operator state, and query metadata. With an idempotent sink: exactly-once;
    without it: at-least-once with duplicates.
    \item \textbf{[I] When is \texttt{availableNow} the right trigger?} (Ch.~13) Whenever
    freshness is minutes-to-hours: streaming sources at batch economics---process the backlog,
    stop, pay nothing until next schedule.
    \item \textbf{[F] What does a watermark actually do?} (Ch.~14) Defines the lateness
    horizon after which events drop, and authorizes eviction of older state---unbounded memory
    becomes a bounded budget.
    \item \textbf{[I] Why is \texttt{dropDuplicates} dangerous without a watermark?} (Ch.~14)
    Deduplication remembers every key ever seen; without eviction, state grows until the
    executors die.
    \item \textbf{[A] Argue exactly-once for a pipeline, end to end.} (Ch.~14) Checkpointed
    offsets, deterministically derived and checkpointed state, idempotent sink (keyed MERGE or
    Delta replace): a crash replays a batch into a sink that ignores replays.
    \item \textbf{[A] A stakeholder wants ``no late data dropped, ever.'' Response?} (Ch.~14)
    Then state is unbounded and windows never close; the honest design is a long watermark
    plus a batch reconciliation path for the long tail.
    \item \textbf{[I] Why a schema registry for streams?} (Ch.~15) Schemas become versioned
    contracts with compatibility rules; a breaking producer change fails at registration
    instead of corrupting every consumer at 03:00.
    \item \textbf{[I] How many partitions should a topic have?} (Ch.~15) Peak rate divided by
    per-partition sustainable throughput, times headroom---measured, then rounded up, because
    partitions only ever get added.
    \item \textbf{[A] Why \texttt{foreachBatch} instead of a native MERGE sink?} (Ch.~16)
    Streaming DataFrames do not support MERGE directly; \texttt{foreachBatch} exposes each
    micro-batch as a batch DataFrame while the checkpoint preserves exactly-once across the
    hand-off.
    \item \textbf{[A] When is continuous processing justified over 30-second micro-batches?}
    (Ch.~16) When the measured business cost of 30 seconds of staleness exceeds the
    premium---rare. Most ``real-time'' requirements are micro-batch requirements.
\end{itemize}

\section{ML and MLOps (Chapters 17--20)}
\index{interview questions!MLOps}%
\begin{itemize}
    \item \textbf{[F] What does autologging not capture that you must add?} (Ch.~17) Your
    context: dataset identity/version, business owner, feature-set version, approval
    linkage---the metadata that makes runs governable.
    \item \textbf{[I] Alias versus version in production code?} (Ch.~17) Always the alias:
    promote and roll back by moving a pointer, never by redeploying version-pinned code.
    \item \textbf{[I] Reproduce a nine-month-old run. What do you need?} (Ch.~17) Git SHA,
    dataset snapshot (or Delta version), parameters, and environment---all four logged.
    \item \textbf{[I] Define point-in-time correctness and its violation cost.} (Ch.~18)
    Every feature for a label at time $t$ computed only from data before $t$; violate it and
    the model trains on the future---inflated offline metrics, useless production behavior.
    \item \textbf{[A] How does a feature store prevent training-serving skew?} (Ch.~18) One
    governed definition feeds both paths: timestamp lookups for training, latest-value lookups
    for serving. Skew requires two definitions; the store has one.
    \item \textbf{[A] A feature job failed silently for three days. What do users see?}
    (Ch.~18) Stale-but-plausible predictions---the dangerous failure. Defense: freshness gates
    in scoring and drift monitoring on feature distributions.
    \item \textbf{[I] Why keep chunks in Delta if the vector index serves queries?} (Ch.~19)
    Delta is the governed, versioned, rebuildable source of truth; the index is a derived
    projection you rebuild when embeddings change or sync corrupts.
    \item \textbf{[I] How do you evaluate a RAG system?} (Ch.~19) A fixed golden set;
    retrieval metrics and LLM-judged generation metrics logged per run; changes gated on the
    metrics in CI.
    \item \textbf{[I] Serverless endpoint versus batch scoring---how do you choose?} (Ch.~20)
    By consumption pattern: interactive apps need endpoints; bulk periodic decisions need
    batch scoring. Many models need both, from the same registry version.
    \item \textbf{[A] Why must promotion stay human-gated?} (Ch.~20) Automated retraining
    proves fit to recent data, not business acceptability---fairness, stability, and threshold
    effects need judgment. Automate up to the gate; gate the alias move.
\end{itemize}

\section{Delivery and CI/CD (Chapter 23)}
\index{interview questions!CI/CD}%
\begin{itemize}
    \item \textbf{[I] Asset Bundles versus Terraform---what goes where?} (Ch.~23) Terraform
    for the slow perimeter (workspaces, principals, metastores); bundles for the fast
    application layer (jobs, pipelines, dashboards). Both in Git, both through review.
    \item \textbf{[F] What does \texttt{bundle validate} actually catch?} (Ch.~23) Schema
    errors, missing variables, bad references, and permission problems---before any cloud
    change. It is the compile step for your platform.
    \item \textbf{[I] Why remote, locked Terraform state?} (Ch.~23) State records what exists;
    local state forks under team/CI use, and a forked state proposes recreating---or
    deleting---your production estate.
    \item \textbf{[A] How do you keep humans out of prod without slowing hotfixes?} (Ch.~23)
    The same pipeline with an expedited approval gate: the hotfix is a PR with a faster review
    SLA, not a different process.
    \item \textbf{[I] How do you attribute a cost spike to a change?} (Ch.~23) Join the
    billing window to deployment history---bundle deploys and Terraform applies are logged, so
    the cost curve diffs against the release curve.
\end{itemize}

\section{Capstone and Architecture Defense (Chapter 24)}
\index{interview questions!architecture defense}%
\begin{itemize}
    \item \textbf{[F] Walk me through your platform end to end.} (Ch.~24) Ninety seconds:
    sources $\rightarrow$ ingestion $\rightarrow$ medallion $\rightarrow$ serving/ML
    $\rightarrow$ governance $\rightarrow$ delivery, one number per arrow. Five minutes means
    you do not yet own it.
    \item \textbf{[A] What would you rebuild differently at 10x scale?} (Ch.~24) Have the real
    answer: partition/parallelism ceilings, orchestration fan-in, the cost model. ``Nothing''
    fails the question.
    \item \textbf{[A] Your most expensive mistake in this build?} (Ch.~24) Name one---the
    untested rollback, the unmonitored endpoint, the leaked token---quantify it, and show the
    mechanism that now prevents it.
    \item \textbf{[I] How does your platform prove data correctness to a skeptic?} (Ch.~24)
    Reconciliation arithmetic per transition, expectation metrics from the event log, replay
    drills, and the audit trail---evidence, not assurances.
    \item \textbf{[A] Why Databricks over a point-solution stack?} (Ch.~24) One governed
    storage and catalog under batch, streaming, SQL, and ML---versus four governance models
    and four copies. Know the honest counter-costs too (DBU economics, Spark operational
    depth).
\end{itemize}

\begin{tipbox}
Mock-interview protocol: draw three questions at random across domains, answer aloud in
ninety seconds each, and have the interviewer probe one level deeper than your answer. Any
question you survive only by reciting---rather than reasoning from a lab you built---goes on
the rebuild list from Appendix~E.
\end{tipbox}
